\section{Method}
\label{sec:method}

\subsection{Setup}
Let $\entity$ index candidate entities, each observed $\nobs_{\entity}$ times, and let $\stat$
be the statistic the screen scores with.

\begin{definition}[Selection operator]
\label{def:selection}
$\stat$ is a selection operator if it returns a function of the largest of several noisy
estimates --- a maximum, a top-$k$ mean, an upper quantile, an enrichment score, or a
best-of-$N$ choice.
\end{definition}

\begin{proposition}[Count-dependent inflation]
\label{prop:inflation}
For a selection operator $\stat$ applied to observations carrying no effect,
$\Ex[\stat \mid \nobs]$ depends on $\nobs$. Scores computed at different $\nobs$ are therefore
not comparable.
\end{proposition}

\subsection{Calibration}
\label{sec:calibration}

Fit the null by resampling observations \emph{known to carry no effect}, applying the same
$\stat$, over a grid of counts spanning the real range. Then
\begin{equation}
\label{eq:z}
  \calibration
\end{equation}
and rank on $\zof{\entity}$, never on $\score$.

\subsection{The direction is measured, not assumed}
\label{sec:direction}

\todo[inline]{The section reviewers will attack. Make the two-slope evidence explicit: it is
the argument that a hard-coded correction is wrong for half the screens it meets.}

\subsection{The ten stages}
\label{sec:stages}

Calibration is \stage{1} of ten. The others exist because skipping each one cost something
measured; the full account is in the repository's methodology document.
